\documentclass[10pt,a4paper]{amsart}
\usepackage[margin=19mm]{geometry}
\usepackage{amsmath,amssymb,mathtools}
\usepackage[scr=rsfs,cal=euler]{mathalfa}
\usepackage{xfrac}
\usepackage[unicode,hidelinks,bookmarks=false]{hyperref}
\newcommand{\ie}{i.e.\ }
\newcommand{\cf}{cf.\ }
\newcommand{\resp}{respectively\ }
\newcommand{\Subsection}{\subsection}
\providecommand{\nobreakdash}{\nobreak\hbox{-}}
\setlength{\emergencystretch}{2em}
\multlinegap0pt
\usepackage{amssymb}
\usepackage{algorithm}\usepackage{algpseudocode}
\usepackage{mathtools}
\newtheorem{lemma}{Lemma}
\newtheorem{theorem}{Theorem}
\newcommand{\FF}{\mathbb{F}}
\DeclareMathOperator{\Hull}{Hull}
\DeclareMathOperator{\rank}{rank}
\title{The Closure of LCD-to-GI Reductions via Generalized Inner Products}
\author{Keita Ishizuka}
\date{Source: arXiv:2605.23120v1 (CC BY 4.0)}
\begin{document}
\maketitle

\noindent\emph{Source and licence.} The Closure of LCD-to-GI Reductions via Generalized Inner Products. Version arXiv:2605.23120v1 (CC BY 4.0), distributed by arXiv under the Creative Commons Attribution 4.0 International licence, http://creativecommons.org/licenses/by/4.0/, as recorded in the arXiv licence metadata for this exact version and shown on the abstract page https://arxiv.org/abs/2605.23120v1. Full licence notice: \texttt{licenses/arxiv-2605.23120-CC-BY-4.0.txt}.\par\smallskip

\noindent\textit{Editorial scope.} Counted unit: the article's theorem thm:generalized\_massey (source0.tex lines 203--208). The article's complete original proof of it is delivered with it (source0.tex lines 210--225, one \texttt{proof} environment of 16 lines). No mathematics is written, repaired or re-derived: the statement and the proof are the article's own lines, in the article's own notation, and no retained line is substituted or re-flowed.  Retained uncounted as the source-local context the proof needs: lines 195--198.  Nothing was omitted from any retained span, so the package carries no omission record.  No retained line contains a citation, so the wrapper carries no bibliography and no bibliography entry.  No retained line contains a figure environment, an image-inclusion command or drawing code, so no image is delivered and the package declares no asset dependency.  Editorial formatting only, in addition: an AMS \texttt{amsart} preamble that restates the article's own declarations and macros used by the retained lines, namely the environments \texttt{lemma}, \texttt{theorem} on separate counters, together with the standard packages those lines need.  The retained environments print in this document as Lemma 1, Theorem 1 in order of appearance, whereas the article prints the same items with the numbering of its own full document.  The complete native source of the article as distributed in the arXiv e-print for this exact version is bundled unchanged as \texttt{sources/201186/source0.tex}.
\begin{lemma}
\label{lem:hull_dimension}
Let $C$ be an $[n,k]$ code with generator matrix $G$, and let $M$ be a structure matrix. Then $\dim \Hull_M(C) = k - \rank(GMG^T)$.
\end{lemma}

\begin{theorem}[Generalized Massey Theorem]
\label{thm:generalized_massey}
Let $C$ be an $[n,k]$ linear code with generator matrix $G$, and let $M$ be a nondegenerate symmetric structure matrix. Then $\Hull_M(C) = \{0\}$ if and only if $GMG^T$ is nonsingular. Moreover, if $\Hull_M(C) = \{0\}$, then
$$\Pi_{C,M} = MG^T(GMG^T)^{-1}G$$
is the orthogonal projector from $\FF_q^n$ onto $C$ with respect to $M$, satisfying $\Pi_{C,M}^2 = \Pi_{C,M}$, $\text{Im}(\Pi_{C,M}) = C$, and for all $v, w \in \FF_q^n$ with $v\Pi_{C,M} = v$ and $w\Pi_{C,M} = 0$, we have $vMw^T = 0$.
\end{theorem}

\begin{proof}
By Lemma~\ref{lem:hull_dimension}, $\dim \Hull_M(C) = k - \rank(GMG^T)$. Thus $\Hull_M(C) = \{0\}$ if and only if $\rank(GMG^T) = k$, i.e., $GMG^T$ is nonsingular.
For the projector formula, we verify that $\Pi_{C,M} = MG^T(GMG^T)^{-1}G$ satisfies the three defining properties. Idempotence $\Pi_{C,M}^2 = \Pi_{C,M}$ follows by direct computation.

For the image, let $v = cG \in C$ where $c \in \mathbb{F}_q^k$.
Then
\begin{align*}
v\Pi_{C,M}
&= cG \cdot MG^T(GMG^T)^{-1}G\\
&= c(GMG^T)(GMG^T)^{-1}G\\
&= cG = v,
\end{align*}
so $v$ lies in the image of $\Pi_{C,M}$. Conversely, any vector in the image has the form $w\Pi_{C,M} = wMG^T(GMG^T)^{-1}G$ for some row vector $w \in \mathbb{F}_q^n$; setting $c = wMG^T(GMG^T)^{-1} \in \mathbb{F}_q^k$ gives $w\Pi_{C,M} = cG \in C$. Hence the image equals $C$.

For orthogonality, let $v, w \in \mathbb{F}_q^n$ satisfy $v\Pi_{C,M} = v$ and $w\Pi_{C,M} = 0$. Then $v = cG$ for some $c \in \mathbb{F}_q^k$, and from $wMG^T(GMG^T)^{-1}G = 0$ we obtain $wMG^T = 0$ since $(GMG^T)^{-1}$ is invertible and $G$ has full row rank. Therefore $vMw^T = cGMw^T = c(wMG^T)^T = 0$, as required.
\end{proof}

\end{document}
